\documentclass[11pt,a4paper]{article}

\usepackage[margin=2.6cm]{geometry}
\usepackage{amsmath,amssymb,amsthm}
\usepackage{hyperref}
\usepackage{booktabs}
\usepackage{enumitem}
\usepackage{natbib}
\bibliographystyle{plainnat}

\newtheorem{theorem}{Theorem}[section]
\newtheorem{proposition}[theorem]{Proposition}
\newtheorem{lemma}[theorem]{Lemma}
\newtheorem{corollary}[theorem]{Corollary}
\theoremstyle{remark}
\newtheorem{remark}[theorem]{Remark}

\newcommand{\Q}{\mathbb{Q}}
\newcommand{\Z}{\mathbb{Z}}
\newcommand{\N}{\mathbb{N}}
\newcommand{\fract}[1]{\left\{#1\right\}}
\newcommand{\lean}[1]{\texttt{#1}}
\newcommand{\Nu}{\mu}

\title{Zero-Entropy Cages and Discrete Fluid for the Mahler $\tfrac32$ Map}
\author{Su MingKai}
\date{2026}

\begin{document}

\begin{abstract}
We study two complementary obstructions for the Mahler map $T x=\fract{3x/2}$: interval
\emph{cages} that admit no long orbit, and a \emph{discrete fluid} picture in which mass is
transported by the piecewise-affine branches.  The exact finite content is machine-checked in
Lean 4 with Mathlib: the integer-vanishing engine, an exact four-term cage for the
\emph{deterministic} even branch on $[2/7,5/7)$, the two-branch joint relation and its genuine
periodic survivors, the exact three-cycle $4/19\mapsto 6/19\mapsto 9/19\mapsto 4/19$ of the
$\bmod\,\tfrac12$ cup map, piecewise expansion by $\tfrac32$, and the $2{:}1$ mass split.  The
``all real numbers are excluded'' statement is isolated as a human-level proposition; the
analytic and spectral numbers sometimes quoted for it are reported as empirical and explicitly
not asserted as theorems.
\end{abstract}

\tableofcontents

\section{Introduction}

The base $\tfrac32$ expansion $\fract{\xi(3/2)^n}$ is the subject of \citep{strauch1997,
lagarias2003,andrieu2026,cassaigne2026}, with the arithmetic side developed in
\citep{stephan2026subspace,stephan2026minimal,stephan2026third,cook2026,calegari2025,elmv}.
The two obstructions treated here are elementary but exact, and they are the part that survives
formalisation without numerical analysis.

\paragraph{Scope and honesty.}
Section~\ref{sec:cages}--\ref{sec:fluid} contain machine-checked results.  The quantitative
claims about chaos onset, spectral gaps, Hausdorff dimension and the cosine-product bound are
\emph{empirical}: they are reported in Section~\ref{sec:empirical} with explicit warnings, and the
accompanying script recomputes the finite certificates but does not turn those numbers into
theorems.  The ``$8/57$'' constant is identified as the exact preimage $T_0^{-1}(4/19)$ of the cup
cycle, not as a critical emptying threshold (Remark~\ref{rem:857}).

\section{Zero-entropy cages}
\label{sec:cages}

\subsection{The engine}

\begin{lemma}[Integer exclusion]
\label{lem:engine2}
If $2^{N}\mid M$ and $|M|<2^{N}$ for $M\in\Z$, then $M=0$.
\end{lemma}

\begin{proof}
$M=2^N k$ gives $2^N|k|<2^N$, hence $|k|<1$ and $k=0$.
Machine-checked: \lean{Mahler32.ZeroEntropyCages.exclusion\_integer}.
\end{proof}

Lemma~\ref{lem:engine2} is the arithmetic half of every exclusion argument: an eventually
periodic orbit produces an integer divisible by all powers of $2$, and the Archimedean bound
kills it.

\subsection{The exact deterministic $[2/7,5/7)$ cage}

Write $Tq$ for the Lean implementation of $T$ on $\Q$,
$Tq(x)=\tfrac{3x}{2}$ if $\tfrac{3x}{2}<1$ and $\tfrac{3x}{2}-1$ otherwise.  The two monotone
branches are $T_0 x=\{3x/2\}$ (equal to $Tq$) and $T_1 x=\{3x/2+1/2\}$; a real Mahler orbit may
switch branches at every step, so the honest object is the \emph{joint relation}
$\mathrm{Step}(x,y)\iff y=T_0x\ \text{or}\ y=T_1x$.  The theorem below concerns $Tq$ alone and
must not be read as an exclusion statement for $\mathrm{Step}$.

\begin{theorem}[No four-term orbit for $Tq$]
\label{thm:cage}
There is no $x\in\Q$ with $x,Tq x,Tq^{2}x,Tq^{3}x$ all in $[2/7,5/7)$; equivalently, if the
first three lie in the interval then the fourth does not.
\end{theorem}

\begin{proof}
One tracks the exact rational images.  From $x\in[2/7,5/7)$ the wrap branch of the first step is
impossible (it would force $x\ge 6/7$), so $x<10/21$ and $Tq x=3x/2$.  The same argument at the
second step forces $Tq^{2}x=9x/4$ and $x<20/63$.  For the third step, $z=9x/4\in[9/14,5/7)$; if
$\tfrac{3z}{2}<1$ then $Tq z=\tfrac{27x}{8}\ge\tfrac{27}{28}>\tfrac57$, and otherwise
$Tq z=\tfrac{27x}{8}-1<\tfrac1{14}<\tfrac27$.  Either way $Tq^{3}x\notin[2/7,5/7)$.
Machine-checked: \lean{Mahler32.ZeroEntropyCages.cage\_two\_sevenths}.
\end{proof}

\begin{remark}[The single-branch theorem does not exclude joint orbits]
\label{rem:joint}
Theorem~\ref{thm:cage} says that the \emph{deterministic} map $Tq$ has no point with four
consecutive hits in $[2/7,5/7)$, so its survivor set is empty.  It does \emph{not} exclude orbits
of the two-branch relation $\mathrm{Step}$.  Indeed
\[
  \tfrac37 \xrightarrow{T_0} \tfrac9{14} \xrightarrow{T_1} \tfrac{13}{28}
  \xrightarrow{T_0} \tfrac{39}{56} \xrightarrow{T_1} \tfrac{61}{112}
  \xrightarrow{T_1} \tfrac{71}{224}
\]
is a six-term $\mathrm{Step}$-orbit inside $[2/7,5/7)$ (machine-checked:
\lean{joint\_orbit\_two\_sevenths}), and $2/5\leftrightarrow 3/5$ and
$4/13\mapsto6/13\mapsto9/13\mapsto7/13\mapsto4/13$ are genuine $\mathrm{Step}$-cycles inside the
same interval (\lean{joint\_cycle\_two\_fifths}, \lean{joint\_cycle\_four\_thirteenths}).  The
survivor set of $\mathrm{Step}$ on $[2/7,5/7)$ is therefore \emph{non-empty}; the script computes
it as an infinite Cantor-like set of sub-exponential cylinder growth (zero topological entropy),
not a finite union of cycles.
\end{remark}

\begin{proposition}[Zero-entropy exclusion, human-level]
\label{prop:allreal}
Let $U\subset[0,1)$ be a finite union of rational intervals whose survivor set
$\bigcap_{n\ge0}T^{-n}U$ is contained in $\{0\}$.  Then no $\xi>0$ satisfies
$\{\xi(3/2)^n\}\in U$ for all $n$.
\end{proposition}

\begin{proof}[Proof sketch]
If a survivor existed with $\xi>0$, zero entropy of the coding would make its $2$-adic expansion
eventually periodic; clearing denominators produces an integer divisible by $2^{N}$ for every
$N$, and Lemma~\ref{lem:engine2} forces it to vanish, i.e.\ $\xi=0$.
\end{proof}

\begin{remark}
Proposition~\ref{prop:allreal} is \emph{not} formalised.  Its second sentence (zero entropy
$\Rightarrow$ eventually periodic) and the clearing step are human-level; only the integer engine
(Lemma~\ref{lem:engine2}) and the concrete cage (Theorem~\ref{thm:cage}) are machine-checked.  We
resist the temptation to present the general statement as a Lean theorem.
\end{remark}

\subsection{Width bookkeeping}

\begin{proposition}
The single interval $[2/7,5/7)$ has width $3/7$; the holed union
$[1/7,3/7)\cup[4/7,6/7)$ has total width $4/7>1/2$.
\end{proposition}

\begin{proof}
Direct arithmetic.  Machine-checked: \lean{width\_three\_sevenths},
\lean{width\_four\_sevenths}, \lean{four\_sevenths\_gt\_half}.
\end{proof}

\section{Chaos onset and the holed $4/7$ cage}
\label{sec:onset}

For a barrier set $U$ the honest object is the \emph{joint survivor sequence} of the two-branch
relation,
\[
  V_0=U,\qquad V_{n+1}=U\cap\bigl(T_0^{-1}V_n\cup T_1^{-1}V_n\bigr),
\]
which is a finite union of rational half-open intervals for every $n$.  The script
\texttt{scripts/generate\_certificates.py} computes $V_n$ exactly, together with a finite-depth
rational Markov partition, its SCC/branching structure, and the exact rational periodic
$\mathrm{Step}$-cycles inside $U$.  The results replace the earlier (and incorrect) ``single-branch
emptying depth'' reading.

\begin{center}
\begin{tabular}{@{}lccc@{}}
\toprule
Barrier set $U$ & $T_0$-alone depth & joint survivor & structure\\
\midrule
$[2/7,5/7)$ & $3$ & infinite & Cantor-like, cubically growing cylinder count\\
$[1/7,3/7)\cup[4/7,6/7)$ & $5$ & infinite & Cantor-like, sub-exponential growth\\
$[28/97,69/97)$ & $3$ & infinite & exactly $2k$ components at depth $k$\\
$[5/24,11/24)\cup[13/24,19/24)$ & $3$ & finite & $\{4/15,2/5,3/5,11/15\}$, one $2$-cycle\\
$[1/5,7/15)\cup[8/15,4/5)$ & $4$ & infinite & linearly growing component count\\
\bottomrule
\end{tabular}
\end{center}

The single-branch depths (iterating $T_0$ alone) are listed only to emphasise that they are
\emph{not} joint emptying depths: every non-trivial barrier keeps a non-empty joint survivor set,
and the ``depth $3$'' of $[2/7,5/7)$ says nothing about the relation $\mathrm{Step}$
(Remark~\ref{rem:joint}).  The survivor cylinder count grows sub-exponentially in each infinite
row (zero topological entropy), but it is infinite, so ``finitely many cycles'' would be false; the
one barrier whose joint survivor set is genuinely finite is the holed cage
$[5/24,11/24)\cup[13/24,19/24)$, whose survivor set is exactly $\{4/15,2/5,3/5,11/15\}$ with the
$2$-cycle $2/5\leftrightarrow3/5$ and the transients $4/15\mapsto2/5$, $11/15\mapsto3/5$.  The Lean
module \lean{Mahler32.ZeroEntropyCages} proves the existence of these cycles and transients
(\lean{holed\_cage\_two\_cycle}, \lean{holed\_cage\_transients}) and the joint orbit witnesses; the
full survivor-count and SCC analysis is the script certificate.

\begin{remark}[The constant $8/57$]
\label{rem:857}
The working notes call $\delta^\ast=8/57$ the critical value of the left-hole family
$[\delta,\tfrac12)$.  Exact computation identifies the arithmetic origin of the constant:
\[
  8/57 \;=\; \tfrac{2}{3}\cdot\tfrac{4}{19} \;=\; T_0^{-1}(4/19),
\]
the exact preimage, under the even branch, of the cup $3$-cycle point $4/19\mapsto6/19\mapsto9/19$.
Under the cup map $F$ --- the restriction of the joint relation to $[0,\tfrac12)$, i.e.\
$F(x)=T_0(x)$ on $[0,\tfrac13)$ and $F(x)=T_1(x)$ on $[\tfrac13,\tfrac12)$ --- the family
$[\delta,\tfrac12)$ has three exact regimes, verified in exact rational arithmetic:
\begin{itemize}[leftmargin=1.4em]
  \item $\delta>4/19$: $[\delta,\tfrac12)$ empties in finitely many steps;
  \item $8/57<\delta\le 4/19$: it does \emph{not} empty --- the $3$-cycle
        $4/19\mapsto6/19\mapsto9/19\mapsto4/19$ survives --- but the survivor set is finite, hence of
        zero topological entropy;
  \item $\delta<8/57$: the fold $T_1\circ T_0^3$ sends $[8/81,8/57)$ into $[0,4/19)$ (indeed
        $T_1(T_0^3(8/57))=4/19$ and $T_1(T_0^3(8/81))=0$), creating a feedback loop and positive
        topological entropy.
\end{itemize}
Thus $8/57$ is a genuine threshold \emph{for the cup map}, but the mechanism is the fold
$T_1\circ T_0^3$ and the resulting feedback, not a ``single-branch emptying'' phenomenon.
\end{remark}

\section{Discrete fluid on the Mahler cup}
\label{sec:fluid}

The ``cup'' map is $Tcup(x)=(3x/2)\bmod\tfrac12$ on $[0,\tfrac12)$.

\begin{proposition}[Self-map]
If $0\le x<\tfrac12$ then $0\le Tcup(x)<\tfrac12$.
\end{proposition}
\begin{proof}
On the non-wrapping branch $\tfrac{3x}{2}\in[0,\tfrac12)$; on the wrapping branch
$\tfrac{3x}{2}-\tfrac12\in[0,\tfrac14)$.  Machine-checked:
\lean{Mahler32.DiscreteFluidCertificates.Tcup\_mem}.
\end{proof}

\begin{theorem}[Exact three-cycle]
\label{thm:cycle}
$Tcup(4/19)=6/19$, $Tcup(6/19)=9/19$, $Tcup(9/19)=4/19$; the period is exactly three.
\end{theorem}
\begin{proof}
Direct rational evaluation; the return values avoid the wrap until the last step, where
$\tfrac32\cdot\tfrac9{19}-\tfrac12=\tfrac{4}{19}$.
Machine-checked: \lean{cup\_cycle}, \lean{cup\_period\_three}.
\end{proof}

\begin{remark}
Theorem~\ref{thm:cycle} is the reason an ``all positive reals excluded'' statement must respect
the branch structure: the cup map genuinely has a positive periodic survivor.  The cages of
Section~\ref{sec:cages} refer to $T$ on $[0,1)$, not to $Tcup$.
\end{remark}

\begin{proposition}[Expansion and mass split]
On a fixed branch, $Tcup$ stretches distances by exactly $3/2$:
$|Tcup(x)-Tcup(y)|=\tfrac32|x-y|$.  Moreover
$\tfrac12-\tfrac13=2(\tfrac14-\tfrac16)$, and $x\mapsto \tfrac{2x}{3}$ maps
$[1/4,1/2]$ onto $[1/6,1/3]$ with the lower half onto $[1/6,1/4]$; this is the arithmetic content
of $\Nu([1/6,1/4])=\Nu([1/3,1/2])$ for a measure carried by the transport.
\end{proposition}
\begin{proof}
Subtract the two linear branches.  Machine-checked: \lean{Tcup\_dist}, \lean{two\_one\_split},
\lean{scale\_two\_thirds}.
\end{proof}

\section{The $9/4$ model and the Euler triangle}
\label{sec:ninefour}

The linear model $4g_{k+1}=9g_k+d_k$ with $d_k\in\{2,3\}$ is the affine shadow of the cup map
near the three-cycle; the $180^\circ$ mirror symmetry quoted in the notes is the involution
$(\varepsilon,r)\mapsto(1-\varepsilon,4^{k}-1-r)$ on the digit string.  We record it here as
motivation only: the exact statements formalised are Theorem~\ref{thm:cycle} and the two
propositions of Section~\ref{sec:fluid}.  The ``Euler-anchor / triangle'' interpretation, and the
cosine-product bound $\le 0.7409$ sometimes attached to it, are numerical observations
(Section~\ref{sec:empirical}), not Lean theorems.

\section{Empirical observations and open problems}
\label{sec:empirical}

The following numbers appear in the working notes and are \emph{reported}, not proved:
\begin{itemize}[leftmargin=1.4em]
  \item chaos onset $L_{\text{chaos}}\approx 0.428705$;
  \item mixing spectral gap $\approx 0.2898$;
  \item Hausdorff dimension migration $\approx 0.8548$;
  \item a three-step cosine-product bound $\le 0.7409$.
\end{itemize}
They require the Fourier/transfer-operator analysis of the map, which is outside the formalised
core.  The constant $8/57$ is the exact preimage $T_0^{-1}(4/19)$ of the cup cycle
(Remark~\ref{rem:857}), not a threshold.  We list these as open: any of them could be made a
theorem with a suitable analytic development, but none is asserted here.

\section{Finite-roof mirror-block exclusion and carry-free base-9 conservation}
\label{sec:model94}

The integer model $4g_{k+1}=9g_k+d_k$, $d_k\in\{2,3\}$, keeps only $g_k\bmod 4\in\{1,2\}$ (a
residue $1$ adds $3$, a residue $2$ adds $2$).  The module \lean{Mahler32.Model94Obstructions}
formalises seven groups of results with no \lean{sorry}, no \lean{native\_decide} and no custom
axiom; the kernel audit prints only \lean{propext}, \lean{Quot.sound}, \lean{Classical.choice}.

\begin{itemize}[leftmargin=1.4em]
  \item \textbf{Mirror conservation} (\lean{mirror\_conservation\_94}).  For $g_0<4^N$ and $k\le N$,
        if $g_0$ survives $k$ steps to $g_k$, its mirror twin $4^N-1-g_0$ survives to
        $\tilde g_k$ with the exact identity $g_k+\tilde g_k=9^k4^{N-k}-1$.
  \item \textbf{Even-rung annihilation} (\lean{even\_rung\_mirror\_annihilation}).  For even
        $N=2m$, no $g_0$ has both $g_0$ and its mirror twin surviving $N+2$ steps.
  \item \textbf{Forbidden sums} (\lean{forbidden\_sum\_c2\_mod4}).  If $c\equiv2\pmod4$ and
        $g_0+g_0'=c\,4^k-1$, the two cannot both survive $k+1$ steps.
  \item \textbf{Odd-rung Klein-4} (\lean{odd\_rung\_klein4\_no\_four}).  For \emph{every} odd
        $N\ge3$, $g<4^{N-1}$ and $\delta\in\{1,2,3\}$, the four sibling/mirror points
        $g,\ g+\delta4^{N-1},\ 4^N-1-g,\ 4^N-1-(g+\delta4^{N-1})$ cannot all survive $N+1$
        steps; at most three survive.  The proof is uniform in $N$: mirror conservation at level
        $N$ forces the two $N$-th iterates of each mirror pair to be $\equiv2\pmod4$, and the
        difference shadow gives $b_N-a_N=\operatorname{stepDiff}(\delta9^{N-1})$, whose residue
        modulo $4$ is $2$ (for $\delta=1$), undefined (for $\delta=2$) or $3$ (for $\delta=3$)
        when $N$ is odd --- never $0$.  The $N=3$ \lean{decide} instance
        \lean{odd\_rung\_klein4\_exclusion} is retained as a kernel witness; the script verifies
        odd $N\le9$ directly.
  \item \textbf{Shadow dynamics} (\lean{sum\_shadow\_exact}, \lean{diff\_shadow\_exact}).  The
        single-variable maps $\Delta$ and $S$ reproduce the pair difference and sum exactly, and
        $\Delta$ preserves the order.
  \item \textbf{Carry-free base-9} (\lean{base9\_transducer\_exact},
        \lean{base9\_digit\_sum\_conservation}).  The left-to-right transducer obeys the exact
        per-position identity $4a'+r'=9r+a$; over a whole word,
        $4\cdot(\text{output sum})=8\cdot(\text{remainder sum})+(\text{digit sum})+r_0-r_{\text{final}}$.
  \item \textbf{Finite-roof block and general lifting} (\lean{finite\_roof\_block\_certificate},
        \lean{step94\_add\_four\_pow}, \lean{iter94\_add\_four\_pow},
        \lean{survives94\_add\_four\_pow}).  A \lean{decide} check in the $4^4=256$ window gives
        $8\,|S_{N+2}|\le3\,|S_N|$, breaking the $1/2$ mirror bound.  The general lifting engine
        proves that survival to any fixed horizon $k$ is periodic with period $4^k$: adding
        $4^k t$ shifts the $k$-th iterate by $9^k t$, so the finite window $[0,4^k)$ exhausts all
        survival patterns at horizon $k$.  The roof ratio at larger moduli ($4^6,4^7,4^9$) is a
        script certificate.
\end{itemize}

\paragraph{Lean versus script.}
Groups 1--3 and 5--6 are theorems over all $N$; group 4 is now a uniform-in-odd-$N$ theorem
(\lean{odd\_rung\_klein4\_no\_four}) with the $N=3$ \lean{decide} instance kept as a witness; group 7
combines the $4^4$ kernel \lean{decide} check with the general lifting engine
(\lean{survives94\_add\_four\_pow}), whose wider-modulus companions live in
\texttt{scripts/generate\_certificates.py}.  Nothing here claims the Mahler $3/2$ conjecture or the
$9/4$ conjecture.

\section{Formalisation traceability}

\begin{table}[h]
\centering
\small
\begin{tabular}{@{}lll@{}}
\toprule
Result & Lean declaration & Status\\
\midrule
Integer exclusion & \lean{ZeroEntropyCages.exclusion\_integer} & proved\\
$[2/7,5/7)$ deterministic cage & \lean{ZeroEntropyCages.cage\_two\_sevenths} & proved\\
Joint `Step`-orbit in $[2/7,5/7)$ & \lean{joint\_orbit\_two\_sevenths} & proved\\
Joint `Step`-cycles & \lean{joint\_cycle\_two\_fifths}, \lean{joint\_cycle\_four\_thirteenths} & proved\\
Holed-cage $2$-cycle / transients & \lean{holed\_cage\_two\_cycle}, \lean{holed\_cage\_transients} & proved\\
Widths & \lean{width\_three\_sevenths}, \lean{width\_four\_sevenths} & proved\\
Cup self-map & \lean{DiscreteFluidCertificates.Tcup\_mem} & proved\\
$4/19$ cycle & \lean{cup\_cycle}, \lean{cup\_period\_three} & proved\\
Expansion & \lean{Tcup\_dist} & proved\\
$2{:}1$ split & \lean{two\_one\_split}, \lean{scale\_two\_thirds} & proved\\
All-real exclusion & Proposition~\ref{prop:allreal} & not formalised\\
$8/57$ regimes & $8/57=(2/3)(4/19)$, fold $T_1\circ T_0^3$ & script (Remark~\ref{rem:857})\\
Spectral / dimension / cosine & --- & empirical only\\
Mirror conservation & \lean{Model94Obstructions.mirror\_conservation\_94} & proved\\
Even-rung annihilation & \lean{Model94Obstructions.even\_rung\_mirror\_annihilation} & proved\\
Forbidden sums & \lean{Model94Obstructions.forbidden\_sum\_c2\_mod4} & proved\\
Odd-rung Klein-4 & \lean{Model94Obstructions.odd\_rung\_klein4\_no\_four} & proved (all odd $N\ge3$)\\
Shadow dynamics & \lean{sum\_shadow\_exact}, \lean{diff\_shadow\_exact} & proved\\
Carry-free base-9 & \lean{base9\_transducer\_exact}, \lean{base9\_digit\_sum\_conservation} & proved\\
Finite-window lifting & \lean{survives94\_add\_four\_pow} (period $4^k$) & proved (all $k$)\\
Finite-roof block & \lean{finite\_roof\_block\_certificate} & \lean{decide} ($4^4$)\\
\bottomrule
\end{tabular}
\caption{Theorem-to-Lean map for Paper~2.}
\end{table}

\bibliography{refs}

\end{document}
